\section{Evaluation Plan}
\label{sec:evaluation}

The three thrusts will be evaluated through a combination of
controlled testbed experiments, collaborations with industry and
critical-infrastructure partners, and end-to-end case studies.

\paragraph{Reference testbeds.}
We will build a small number of reference \cps testbeds in our lab,
spanning sectors with distinct threat models: a water-treatment
mini-plant (SCADA + PLCs + HMI), a building-automation deployment
(BACnet + IP-connected controllers), and a medical-device bench
(networked infusion pump and monitoring stack, analyzed in cooperation
with our institutional medical-device program). These testbeds are
small enough to allow ground truth about their configuration and
vulnerabilities to be established by hand, but representative enough
that findings generalize. For each testbed we will maintain a
versioned KB and a seeded set of known vulnerabilities (drawn from
historical CVEs applicable to the specific devices and firmware
versions) against which \sysname's precision and recall can be
measured.

\paragraph{Partner environments.}
We will work with industry partners and critical-infrastructure
organizations to evaluate \sysname against real operational environments.
The evaluation will focus on the KB and Observer Agents (which are
the components deployable against production systems) and on the
ability of our framework to replicate partner environments from the
evidence partners already have. We expect partners to provide
documentation and logs, and---under appropriate safeguards---read-only
observation opportunities. The Red Team Agents will be run only
against the resulting twin.

\paragraph{UCSB campus as a secondary target.}
As an intermediate-scale non-\cps environment, we will also develop
Observer Agents for UCSB internal networks. This gives us a large,
heterogeneous, cooperative target on which we can refine observation
techniques before applying them to more sensitive environments.

\paragraph{Metrics.}
We will report the following metrics across thrusts:
\emph{KB coverage and freshness} (fraction of ground-truth entities
represented, median staleness of their properties);
\emph{observation risk} (measured operational impact on targets, with
zero incidents as a hard requirement);
\emph{twin fidelity} (attack reproducibility and behavioral divergence,
as defined in \S\ref{sec:thrust2});
\emph{analysis yield} (number of confirmed vulnerabilities per
twin-hour of analysis); and
\emph{human efficiency} (findings validated per analyst-hour, and
true-positive rate of escalated findings).

\paragraph{Baselines and comparisons.}
Where possible, we will compare \sysname's findings against the
results of traditional manual penetration tests performed on the same
testbeds. We will also compare against purely component-level
analyses (fuzzing alone, static analysis alone), to measure the
marginal contribution of multi-step planning. Finally, we will
conduct ablation studies to quantify the contribution of each agent
role and of the LLM/classical-planner hybrid.
